\documentclass[11pt,a4paper]{article}

\usepackage[a4paper,top=18mm,bottom=20mm,left=17mm,right=17mm,headheight=15pt]{geometry}
\usepackage{fontspec}
\usepackage{xcolor}
\usepackage{array}
\usepackage{tabularx}
\usepackage{enumitem}
\usepackage{fancyhdr}
\usepackage{hyperref}
\usepackage{needspace}

\setmainfont{Arial}
\setsansfont{Arial}

\definecolor{Primary}{HTML}{235347}
\definecolor{Accent}{HTML}{E98B2A}
\definecolor{Soft}{HTML}{EEF4F1}
\definecolor{LineGray}{HTML}{B7C3BE}
\definecolor{TextGray}{HTML}{4E5A56}

\newcommand{\documentlabel}{Angket Kompetitor GrabFood}

\hypersetup{
  colorlinks=true,
  linkcolor=Primary,
  urlcolor=Primary,
  pdftitle={Angket Kompetitor GrabFood Nasi Gerilya},
  pdfauthor={Instrumen Penelitian Skripsi}
}

\pagestyle{fancy}
\fancyhf{}
\fancyhead[L]{\small\textcolor{Primary}{\documentlabel}}
\fancyhead[R]{\small\textcolor{TextGray}{Versi Juli 2026}}
\fancyfoot[C]{\small\textcolor{TextGray}{Halaman \thepage}}
\renewcommand{\headrulewidth}{0.4pt}
\renewcommand{\footrulewidth}{0pt}

\setlength{\parindent}{0pt}
\setlength{\parskip}{5pt}
\renewcommand{\arraystretch}{1.18}
\hyphenpenalty=10000
\exhyphenpenalty=10000
\emergencystretch=2em
\sloppy
\newcolumntype{Y}{>{\raggedright\arraybackslash}X}

\newcommand{\sectionbar}[1]{%
  \Needspace{4\baselineskip}
  \vspace{5mm}
  \noindent\colorbox{Soft}{%
    \parbox{\dimexpr\linewidth-2\fboxsep\relax}{%
      \large\bfseries\textcolor{Primary}{#1}
    }%
  }\par\vspace{2mm}
}

\newcommand{\choice}{\raisebox{0.15ex}{\fbox{\rule{0pt}{1.05ex}\rule{1.05ex}{0pt}}}\hspace{0.5mm}}

\newcommand{\choiceitem}[1]{\choice #1}

\newcommand{\questionblock}[3]{%
  \Needspace{7\baselineskip}
  \noindent
  \begin{tabularx}{\linewidth}{>{\bfseries\color{Primary}}p{8mm} Y}
    #1 & \textbf{#2}\\
       & {\small\textcolor{TextGray}{#3}}\\
  \end{tabularx}
  \par\vspace{2.4mm}
}

\newcommand{\answerline}{%
  \par\vspace{1mm}
  \noindent\textcolor{LineGray}{\rule{\linewidth}{0.35pt}}\par\vspace{4mm}
}

\begin{document}

\begin{titlepage}
\thispagestyle{empty}
\vspace*{12mm}
{\Huge\bfseries\textcolor{Primary}{Angket Kompetitor}}\par
{\Huge\bfseries\textcolor{Primary}{GrabFood}}\par
\vspace{3mm}
{\Large\textcolor{Accent}{Instrumen Ringkas Informan Pembanding}}\par
\vspace{4mm}
{\large\textcolor{TextGray}{Versi terbaru: 01 Juli 2026}}\par
\vspace{7mm}
\textcolor{Accent}{\rule{0.92\linewidth}{1.4pt}}\par
\vspace{10mm}

\begin{tabularx}{\linewidth}{>{\bfseries\color{Primary}}p{38mm} Y}
Topik penelitian & Strategi pemasaran Rumah Makan Nasi Gerilya pada platform GrabFood untuk meningkatkan penjualan berdasarkan analisis SWOT.\\
Sasaran informan & Kasir, waitress, admin order online, atau karyawan rumah makan Padang/merchant makanan sejenis.\\
Metode pengisian & Pilih jawaban yang paling sesuai dengan pengalaman informan. Untuk pertanyaan terakhir, jawab secara singkat berdasarkan pandangan informan.\\
Kerahasiaan & Jawaban digunakan hanya untuk kepentingan penelitian. Informan tidak diminta membuka data penjualan, margin, atau rahasia usaha.\\
\end{tabularx}

\vfill
\colorbox{Soft}{%
  \parbox{\dimexpr\linewidth-2\fboxsep\relax}{%
    \raggedright
    \textbf{\textcolor{Primary}{Petunjuk singkat:}}
    Beri tanda centang pada pilihan jawaban yang paling sesuai. Jika memilih "Lainnya", tuliskan jawaban pada ruang yang tersedia.
  }%
}
\vspace*{8mm}
\end{titlepage}

\sectionbar{Identitas Singkat}
\begin{tabularx}{\linewidth}{>{\bfseries\color{Primary}}p{35mm} Y}
Kode informan & \textcolor{LineGray}{\rule{\linewidth}{0.35pt}}\\[4mm]
Nama usaha & \textcolor{LineGray}{\rule{\linewidth}{0.35pt}}\\[4mm]
Tanggal & \textcolor{LineGray}{\rule{\linewidth}{0.35pt}}\\[4mm]
\end{tabularx}

\sectionbar{Pertanyaan Angket Kompetitor}

\questionblock{1.}{Posisi Kakak di sini apa?}{\choiceitem{Kasir}\quad \choiceitem{Waitress}\quad \choiceitem{Admin order online}\quad \choiceitem{Lainnya: \textcolor{LineGray}{\rule{28mm}{0.35pt}}}}

\questionblock{2.}{Sudah berapa lama bekerja di sini?}{\choiceitem{\textless{} 6 bulan}\quad \choiceitem{6-12 bulan}\quad \choiceitem{\textgreater{} 1 tahun}}

\questionblock{3.}{Pesanan paling ramai biasanya dari mana?}{\choiceitem{Makan di tempat}\quad \choiceitem{Bungkus}\quad \choiceitem{GrabFood}\quad \choiceitem{GoFood}\quad \choiceitem{ShopeeFood}}

\questionblock{4.}{Jam paling ramai biasanya kapan?}{\choiceitem{Pagi}\quad \choiceitem{Siang}\quad \choiceitem{Sore}\quad \choiceitem{Malam}\quad \choiceitem{Akhir pekan}}

\questionblock{5.}{Menu apa yang paling sering dipesan pelanggan?}{Sebutkan 2-3 menu.}
\answerline

\questionblock{6.}{Menurut Kakak, kenapa pelanggan sering memilih menu itu?}{\choiceitem{Rasa}\quad \choiceitem{Porsi}\quad \choiceitem{Harga}\quad \choiceitem{Cepat siap}\quad \choiceitem{Promo}\quad \choiceitem{Sudah terkenal}}

\questionblock{7.}{Rentang harga menu yang paling sering dibeli pelanggan biasanya berapa?}{\choiceitem{\textless{} Rp20.000}\quad \choiceitem{Rp20.000-Rp35.000}\quad \choiceitem{Rp35.000-Rp50.000}\quad \choiceitem{\textgreater{} Rp50.000}}

\questionblock{8.}{Apakah pelanggan sering memakai promo/voucher aplikasi?}{\choiceitem{Sering}\quad \choiceitem{Kadang-kadang}\quad \choiceitem{Jarang}\quad \choiceitem{Tidak tahu}}

\questionblock{9.}{Kalau pesanan online ramai, kendala yang paling sering muncul apa?}{\choiceitem{Menu habis}\quad \choiceitem{Salah item}\quad \choiceitem{Lama siap}\quad \choiceitem{Catatan pelanggan terlewat}\quad \choiceitem{Driver menunggu}\quad \choiceitem{Lainnya: \textcolor{LineGray}{\rule{20mm}{0.35pt}}}}

\questionblock{10.}{Keluhan pelanggan yang paling sering terdengar biasanya tentang apa?}{\choiceitem{Rasa}\quad \choiceitem{Porsi}\quad \choiceitem{Harga}\quad \choiceitem{Lama}\quad \choiceitem{Pesanan kurang lengkap}\quad \choiceitem{Kemasan}\quad \choiceitem{Jarang ada keluhan}}

\questionblock{11.}{Sebelum pesanan diberikan ke pelanggan/driver, biasanya dicek dulu apa saja?}{\choiceitem{Item pesanan}\quad \choiceitem{Lauk}\quad \choiceitem{Kuah/sambal}\quad \choiceitem{Minuman}\quad \choiceitem{Catatan khusus}\quad \choiceitem{Tidak selalu dicek}}

\questionblock{12.}{Menurut Kakak, hal apa yang paling membuat pelanggan balik lagi ke restoran ini?}{\choiceitem{Rasa}\quad \choiceitem{Harga}\quad \choiceitem{Porsi}\quad \choiceitem{Pelayanan cepat}\quad \choiceitem{Reputasi restoran}\quad \choiceitem{Lokasi}\quad \choiceitem{Promo}}

\sectionbar{Pertanyaan Terbuka}
\Needspace{8\baselineskip}
\textbf{Kalau dibanding rumah makan Padang lain di sekitar sini, menurut Kakak keunggulan restoran ini apa?}
\answerline
\answerline

\end{document}
